\subsection{The relation between the vertical ideal and the standard mesh}

The vertical-ideal class and the standard mesh respond differently
to ramification.  An exact relation between them explains both a
family of vanishings and a failure of splitting for a particular
nonzero class.

\begin{theorem}\label{thm:vertical-standard-relation}
Let $O$ be a mixed-characteristic DVR of residue characteristic $p$,
with uniformizer $\pi$ and $h=v_\pi(p)$.  Let $n=ps$ with $s\geq1$,
and put $A=O[t]/t^n$, $R=A[C_p]$, $X=\sigma-1$, and $Y=X+t^s$.
Let
\[
 y=[(\pi,Y)]-[R],\qquad z=[M_1]-[R],
\]
where the standard family $M_i$ is defined by
\eqref{eq:pure-power-mesh-presentation}.  Then
\begin{equation}\label{eq:vertical-standard-relation}
 h y=s z\quad\text{in }K_0(\Lat_A(C_p)).
\end{equation}
The same equality holds after induction to $S_p$.

If $p\mid h$, then $s z=0$.  In particular, if $n=p$ and
$p\mid h$, every standard-mesh difference $[M_i]-[R]$ vanishes.
If $p\nmid h$ and $p^2\mid n$, then $y$ and its induction are
nonzero elements of order $p$ which are divisible by $p$ in their
respective whole Grothendieck groups.  The cyclic subgroup generated
by either of these particular classes is consequently not a direct
summand.
\end{theorem}

\begin{proof}
The proof takes place over the original DVR.  No passage to a
normalization or a completion is needed.  For $0\leq a\leq h$ put
\[
 J_a=(\pi^a,Y)\subset R.
\]
The monic relation for $Y$ has every lower coefficient divisible by
$p$: modulo $p$ it is $Y^p-t^n=Y^p$.  It follows that $J_a$ is
$A$-free of rank $p$, with basis
\begin{equation}\label{eq:vertical-power-basis}
 \pi^a,Y,Y^2,\ldots,Y^{p-1}.
\end{equation}
Indeed, reducing modulo $\pi^a$ identifies the ambient algebra
with $(A/\pi^aA)[Y]/Y^p$.  The inverse image of the ideal $(Y)$
consists precisely of elements whose constant coefficient in the
monic $Y$-basis is divisible by $\pi^a$.  This proves spanning in
\eqref{eq:vertical-power-basis}; independence follows because
$\pi$ is a nonzerodivisor on $A$.  We have $J_0=R$ and
$J_1=(\pi,Y)$.

Write $H(X)=g(X)/X$, where $g(X)=(1+X)^p-1$.  There is a polynomial
$V(t)\in O[t]$ with constant term $-1$ such that
\begin{equation}\label{eq:vertical-parameter-unit}
 t^{n-s}-H(-t^s)=pV(t).
\end{equation}
For odd $p$, the leading term of $H(-t^s)$ is $t^{n-s}$,
and every other binomial coefficient in $H$ is divisible by $p$.
For $p=2$, the left side is $2(t^s-1)$, giving $V=t^s-1$.
These computations prove the identity and its asserted constant
term at every prime.  Thus $V$ is a unit in $A$, where $t$ is
nilpotent.  Evaluating the group relation at $X=-t^s$ gives
\[
 g(-t^s)=-t^sH(-t^s)=p t^sV(t),
\]
since $t^n=0$.  Consequently
\begin{equation}\label{eq:vertical-power-quotient}
 R/(Y)\cong A/(p t^s)\cong A/(\pi^h t^s).
\end{equation}

For $1\leq a<h$ the sequence
\begin{equation}\label{eq:vertical-power-mesh}
 0\longrightarrow J_a
 \xrightarrow{u\mapsto(u,\pi u)}J_{a-1}\oplus J_{a+1}
 \xrightarrow{(v,w)\mapsto\pi v-w}J_a
 \longrightarrow0
\end{equation}
is exact.  Injectivity of the first map is immediate, and
surjectivity of the second follows from
$\pi J_{a-1}+J_{a+1}=J_a$.  To check its kernel, let
$v\in J_{a-1}$ satisfy $\pi v\in J_{a+1}$.  Modulo $Y$,
this is a colon computation in the ring
\eqref{eq:vertical-power-quotient}.  In $A$ it reads
\[
 \pi b\in\pi^{a+1}A+\pi^h t^sA.
\]
Cancel the nonzerodivisor $\pi$ in $A$ to obtain
\[
 b\in\pi^aA+\pi^{h-1}t^sA=\pi^aA,
\]
because $a\leq h-1$.  Therefore $v\in J_a$, proving exactness
in the middle.  A coefficient section sends the basis vector
$\pi^a$ to $(\pi^{a-1},0)$ and $Y^q$, $1\leq q<p$, to
$(0,-Y^q)$.  Thus \eqref{eq:vertical-power-mesh} is admissible.

Writing $c_a=[J_a]-[R]$, these relations and $c_0=0$ give
$c_a=a c_1=a y$ for $0\leq a\leq h$.  This also holds when
$h=1$, when there is no interior mesh.  In particular
\begin{equation}\label{eq:vertical-power-endpoint}
 [J_h]-[R]=h y,\qquad J_h=(p,Y).
\end{equation}

We identify that endpoint with a standard-mesh lattice.  Put
\[
 B(X,t)=t^{n-s}-H(X).
\]
Polynomial division by the monic linear polynomial $Y=X+t^s$
and \eqref{eq:vertical-parameter-unit} give
\[
 B=pV+YW\quad\text{for some }W\in A[X].
\]
Since $V$ is a unit, $(B,Y)=(p,Y)=J_h$.
The standard lattice $M_s$ has relations
\[
 Xe=t^{n-s}f,\qquad t^se+H(X)f=0.
\]
The assignment
\begin{equation}\label{eq:vertical-standard-isomorphism}
 e\longmapsto B,\qquad f\longmapsto Y
\end{equation}
respects them, because in $R$
\[
 XB=Xt^{n-s}-g(X)=t^{n-s}Y,\qquad
 t^sB+H(X)Y=t^n+g(X)=0.
\]
Its image is $J_h$.  It is an isomorphism: the source basis is
$e,X^jf$ for $0\leq j\leq p-2$, the target basis is
$p,Y,\ldots,Y^{p-1}$, and the last $p-1$ image vectors are
$Y(Y-t^s)^j$.  They give a triangular change of basis with
unit diagonal on the positive powers of $Y$, while the first
image has constant coordinate $V$ in the target basis.
The determinant of the full coefficient matrix is therefore $V$,
a unit.  This proves \eqref{eq:vertical-standard-isomorphism}
as an isomorphism of coefficient-free $R$-modules.

The standard mesh gives $[M_s]-[R]=s z$ by
\eqref{eq:pure-power-mesh-relations}.  Combining this with
\eqref{eq:vertical-power-endpoint} proves
\eqref{eq:vertical-standard-relation}.  Induction is exact and
preserves this equality.

By Theorem~\ref{thm:vertical-ideal-exact-order}, $py=0$ and
both $y$ and its induced class have exact order $p$.
If $p\mid h$, the left side of
\eqref{eq:vertical-standard-relation} is zero, giving $s z=0$.
For $n=p$, one has $s=1$, so $z=0$ and then every standard
mesh difference is zero.

Finally suppose $p\nmid h$ and $p^2\mid n$.  Choose an integer
$b$ with $bh\equiv1\pmod p$.  Since $py=0$, the relation gives
\[
 y=bh y=bs z
    =p\left(b\frac{s}{p}z\right).
\]
The same equality holds after induction, with a nonzero left side
in both cases.  Every homomorphism to $\Z/p$ kills a
$p$-divisible element.  If the cyclic subgroup generated by
that element were a direct summand, its projection, followed by
identification of its generator with $1\in\Z/p$, would contradict
this fact.  This proves the nonsplitting assertion for the specified
classes.
\end{proof}

\begin{remark}
The nonsplitting assertion concerns the subgroup generated by the
vertical-ideal class.  For example, over $\Z_p$ the whole generic
kernel already contains split copies of $\Z/n$.  If $p^2\mid n$,
the relation is $y=(n/p)z$, so the nonzero subgroup generated by
$y$ lies inside such a cyclic group as a $p$-divisible subgroup.
There is no claim here that the whole Grothendieck group has no
other split subgroup of order $p$.
\end{remark}

\begin{corollary}\label{cor:vertical-primary-split}
In the notation of Theorem~\ref{thm:vertical-standard-relation}, write
$n=p^k m$ with $p\nmid m$, and suppose $p\nmid h$.
The classes $w=mz$ and $w_{S_p}=\operatorname{Ind}_{C_p}^{S_p}w$
have exact order $p^k$ and generate subgroups split in the respective
whole Grothendieck groups.  They lie in both the generic and the
$t$-reduction kernels.  The symmetric-group retraction annihilates
the Young span.  The subgroup generated by $y$, or by its induction,
is the order-$p$ subgroup of the corresponding displayed cyclic
group of order $p^k$.
\end{corollary}

\begin{proof}
Let $\eta$ and $\Theta$ be the whole-group detectors of
Theorem~\ref{thm:ramified-vertical-truncations}, normalized to
take $y$ and its induction to $s=n/p$ in $\Z/n$.
Applying them to $hy=sz$, and to its induced equality, gives
\[
 s(\eta(z)-h)=0,\qquad
 s(\Theta(\operatorname{Ind}z)-h)=0\quad\text{in }\Z/n.
\]
Since $n=ps$, both detector values are congruent to $h$ modulo $p$.
Their reductions modulo $p^k$ are therefore units.  The standard
mesh gives $nz=0$, hence $p^k w=0$, and the same relation after
induction.  Their detector values modulo $p^k$ are units because
$p\nmid m$.  Thus both source classes have exact order $p^k$.
Reducing each detector modulo $p^k$ and multiplying by the inverse
of its value on the corresponding class gives a whole-group
retraction, with section sending $1$ to that class.

For the kernel assertions, all standard lattices become regular
after inverting $p$.  Indeed, $H(X)/p$ is the trivial central
idempotent of $A[1/p][C_p]$.  On that factor $X=0$ and $H=p$,
so the standard relations eliminate $f$; on the complementary
factor $H=0$ and $X$ is invertible, so they eliminate $e$.
Both give a free rank-one module over the corresponding factor.
For $0<i<n$, reduction at $t=0$ gives
$O\oplus O[X]/H(X)$.  Multiplication by $X$ identifies the second
summand with the augmentation ideal of $O[C_p]$.  The
coefficient-split augmentation sequence therefore makes the reduced
Grothendieck class equal to $[O[C_p]]$.  This proves both vanishings
for $z,w$ and their inductions.

Finally choose $b$ with $bh\equiv1\pmod p$.  The exact order of
$y$ and the relation give
\[
 y=bhy=bsz=b p^{k-1}w,
\]
and the same equality after induction.  As $p\nmid b$, these are
the unique order-$p$ subgroups of the displayed cyclic groups.
For $k>1$ their nonsplitting in the whole group was proved in
Theorem~\ref{thm:vertical-standard-relation}.
\end{proof}
